\begin{abstract}

    University faculty are increasingly burdened by administrative
    ``task-drift,'' spending an estimated 20--40\% of their working hours
    answering routine inquiries such as module logistics and general
    administrative questions. This overhead detracts from high-value research
    and pedagogical innovation, and creates a ``feedback gap'' in which
    professors lack the bandwidth to analyse individual student performance,
    leaving course development stagnant. Students, in turn, experience a
    ``clarity deficit'': delayed or ambiguous digital communication breeds
    confusion and dissatisfaction. High student-to-faculty ratios compound the
    problem by limiting personalised consultation, ultimately stifling student
    creativity and the development of higher-order thinking skills.

    This thesis presents the \emph{AI-Powered Omniscient Teacher Assistant}, a
    system that leverages large language models to address these gaps along two
    fronts. First, it automates the handling of routine administrative and
    module-logistics inquiries through retrieval-augmented generation grounded
    in course materials, reclaiming faculty time for research and teaching.
    Second, it surfaces analytics on student questions and performance, giving
    instructors actionable insight to close the feedback gap while providing
    students with timely, consistent, and personalised responses at scale.

    % TODO: replace the bracketed placeholders below with your actual
    % evaluation setup and findings once the experiments are complete.
    We evaluated the system on [dataset / course cohort], measuring [response
    accuracy, latency, faculty time saved, and student satisfaction]. The
    results show that [the assistant resolved X\% of routine inquiries without
    human intervention and improved Y over the baseline]. These findings
    demonstrate that an AI teaching assistant can meaningfully reduce
    administrative task-drift, narrow the feedback and clarity gaps, and support
    more personalised, higher-order learning even under high
    student-to-faculty ratios.

\end{abstract}
